% Separate post-CP20 candidate. Full audits required before promotion.
\subsection{The final ramification group in the full wild image}

Throughout this subsection $K$ is complete discretely valued of mixed
characteristic $(0,p)$, $p\ge5$, with algebraically closed residue field,
$e_K=p-1$ and $\mu_p\subset K$. We use the normalized hypothetical
different-character comparison of Theorem~\ref{thm:B-RD49}, and assume
that its common residual inertia is irreducible over $\mathbf F_p$.
Set $q=p^2$, $N=q-1$ and $b_0=N(p-1)$. By
Theorem~\ref{thm:B-RD57}, its remaining wild image is
$W=M_2(\mathbf F_p)$, written additively through $1+pB$.
The common tame good field $L/K$ has degree $N$, and the common
division field $F/L$ has group $W$.

\begin{lemma}[Point fields and the final Cartan plane]
\label{lem:B-RD58-final-Cartan}
\status{Proved; full independent root and companion proof, original-source, and exact-code audits passed.}
For every primitive second-level point $P$, the field $L(P)/L$ is
Galois of degree $q$ and has a single lower and upper ramification
break $b_0$. The last nonzero upper group of $F/L$ is
\[
                         W^{b_0}=Z,                    \tag{B-RD58.1}
\]
where $Z$ is the full nonsplit Cartan plane. Consequently all nonzero
Cartan elements have a single common lower break $c>b_0$, and every
element outside $Z$ has lower break strictly less than $c$.
\end{lemma}
\begin{proof}
Fix a primitive $P$ and let $v$ be its nonzero image in
$M[p^2]/pM[p^2]$. Its stabilizer in $W$ is
\[
                   H_v=\{B:Bv=0\},
\]
a two-dimensional subspace. Since $W$ is abelian, its fixed field
$E=L(P)$ is Galois over $L$, of degree $q$.
Lemma~\ref{lem:B-RD55-valuations} gives
\[
 v_E(T(P))=p-1,\qquad
 v_E(T(Q))=q(p-1)\quad(0\ne Q\in M[p]).
\]
Every nonidentity automorphism of $E/L$ displaces $P$ by a nonzero
first-level point. Formal subtraction has the same valuation as
ordinary coordinate subtraction. Expand $T(P)$ in a uniformizer
of $E$, using residue representatives fixed by the Galois group as
in (B-RD55.4). Its valuation $p-1$ is prime to $p$, so an
automorphism of lower break $b$ changes it in valuation $p-1+b$.
It follows that $b=q(p-1)-(p-1)=b_0$ for every nonidentity
automorphism. A single lower break is also the single upper break.

We now use the quotient compatibility of upper ramification groups,
in exactly the totally ramified, perfect-residue-field setting of
\cite[\S3.1, Proposition~3, p.~128]{Serre1961LocalClosed}. The full norm-polynomial proof and Herbrand passage are also given in
\cite[Proposition~3.7, Theorem~3.10 and Proposition~3.13,
pp.~14--16]{ColmezCorpsLocaux}.
For every $v\ne0$, the image of $W^{b_0}$ in $W/H_v$ is all of
$W/H_v$, and its image is zero at every upper index greater than
$b_0$. Since $\bigcap_{v\ne0}H_v=0$, there is no nonzero upper
group after $b_0$. Put $S=W^{b_0}$. Thus
\[
                 S\longrightarrow\mathbf F_q,
                         B\longmapsto Bv
               \quad\hbox{is surjective for every }v\ne0.       \tag{B-RD58.2}
\]

The determinant on $W$ is $1+p\operatorname{tr}(B)$.
The determinant of the good height-two, dimension-one group over
$L$ is cyclotomic, by the determinant argument already proved in
B18 and B49; its finite tame determinant has become trivial over $L$.
Thus its degree-$p$ quotient field is $E_{\rm cyc}=L(\mu_{p^2})$.
This field has degree $p$: the cyclotomic extension over the
unramified coefficient field has ramification degree $p(p-1)$,
whereas $e_L=(p-1)N$ is prime to $p$.
For $x=\zeta_{p^2}-1$ one has
\[
 v_{E_{\rm cyc}}(x)=N,
 \qquad v_{E_{\rm cyc}}(g x-x)=pN\quad(g\ne1).
\]
The same prime-to-$p$ valuation calculation gives its single break
$pN-N=b_0$. Quotient compatibility therefore shows that
$\operatorname{tr}:S\to\mathbf F_p$ is nonzero.

Choose a tame complement of order $N$ and an eigen-uniformizer
$\varpi$ of $F$. As in B56, its multiplier is a primitive $N$th
root $\eta$; here $\eta=\zeta$, since $p^4\equiv1\pmod N$.
If $c$ is the last lower break, the leading ramification map
\[
 S\hookrightarrow k,\qquad
 g\longmapsto d_g=
    \overline{(g\varpi-\varpi)/\varpi^{c+1}}
\]
is injective and additive, with tame action multiplication by
$\eta^c$. Since trace on $S$ is nonzero and tame-invariant,
averaging supplies a nonzero tame-fixed vector of $S$.
Its nonzero image in $k$ forces $\eta^c=1$.
Injectivity then says that tame conjugation fixes all of $S$.
The centralizer of the irreducible tame residual generator is
precisely $Z=\mathbf F_q$. Thus $S\subset Z$.
Surjectivity in (B-RD58.2) forces $\dim_{\mathbf F_p}S\ge2$,
so $S=Z$.

This is the final lower group, and it becomes trivial immediately
after $c$; hence every nonzero element of $Z$ has break $c$.
The increasing Herbrand function for $F/L$ satisfies
$\varphi(c)=b_0$. Since $S$ is a proper subgroup of $W$,
its slope is strictly less than one on a nonempty interval before
$c$, so $c>b_0$. All other elements disappear before this last
group, proving the final assertion.
\end{proof}

\begin{lemma}[The earlier irreducible anti-Cartan layer]
\label{lem:B-RD58-anti-signature}
\status{Proved; full independent root and companion proof, original-source, and exact-code audits passed.}
Suppose the common residual projective order is $h>2$.
The anti-Cartan plane $U$ has one lower break $b<c$, and
$0<b<b_0$. In an exact Honda-framed coordinate, write the leading
coordinate of a primitive point as $d z\varpi^A$, where
\[
                     A=q(p-1),\qquad z\in\mathbf F_q.
\]
For each $z\ne0$, the first nonzero term of degree prime to $p$
below $qA$ has degree
\[
                     K=A+Nb,                            \tag{B-RD58.3}
\]
and its coefficient is one fixed nonzero multiple of $z$.
Every nonzero term of degree $J$ with $v_p(J)=1$ satisfies
\[
                     J>A+(q-p)b.                        \tag{B-RD58.4}
\]
\end{lemma}
\begin{proof}
The action on $U$ is the irreducible two-dimensional action by
$\lambda^{1-p}$ from B56. Thus every nonzero element of $U$ has
the same lower break $b$, and $b<c$ by the preceding lemma.
Its leading ramification coefficient space
$D_U=\{d_g:g\in U\}\subset k$ is a two-dimensional
$\mathbf F_p$-space stable under multiplication by $\eta^b$.
This scalar has degree two over $\mathbf F_p$, so
\[
                 D_U=d_U\mathbf F_q\quad(d_U\ne0).
\]
In particular $\eta^b\ne1$, so $b\ne b_0$.

Primitive coordinates have valuation $A=q(p-1)$ and nonzero
first-level coordinates have valuation $qA$, by B55. Moreover
$e_F=NqA>qA$, so the fixed-coefficient expansions modulo
$\varpi^{qA+1}$ are characteristic-$p$ expansions. Their coefficients
below $qA$ depend only on the primitive residual class $z$, exactly
as in B56. Every nonzero anti-Cartan matrix is invertible, so
\[
                      v_F(gT(P)-T(P))=qA\quad(g\in U\setminus0).
                                                               \tag{B-RD58.5}
\]
For a monomial of degree $j$, its first change has degree
$j+p^{v_p(j)}b$. The leading degree $A$ has exact $p$-valuation
two, and changes with coefficient $(p-1)d z d_g^q$ in degree
$A+qb$. Every other degree divisible by $q$ changes strictly later.
Write $J$ for the first nonzero degree with $v_p(J)=1$, and $K$
for the first nonzero degree prime to $p$; allow either to be
infinite if absent below $qA$.

If $b>b_0$, then $A+qb>qA$. Any earlier change could only come
from the degree-$J$ and degree-$K$ terms. If their first changes
coincide, their coefficient is a polynomial
$a d_g^p+b' d_g$ with coefficients independent of $g$.
A nonzero polynomial of degree at most $p$ cannot vanish at all
$q$ elements of $D_U$. Therefore (B-RD58.5) forces
$J+pb\ge qA$ and $K+b\ge qA$.
For a nonzero Cartan element of break $c>b$, all these changes
then occur strictly after $qA$, as do those of every degree
divisible by $q$. This contradicts its nonzero first-level
displacement. Hence $b>b_0$ is impossible.

For $b<b_0$, the leading degree changes before $qA$.
There cannot be an even earlier change from $J$ and $K$, by
the same degree-$p$ polynomial argument. Cancellation in degree
$A+qb$ must therefore involve a polynomial
\[
              (p-1)d z X^q+a X^p+b'X
\]
which vanishes on $D_U=d_U\mathbf F_q$.
The monic polynomial with that root space is
$X^q-d_U^{q-1}X$. Comparing coefficients shows that the
$X^p$ coefficient must be zero and the linear coefficient
must cancel the nonzero $X^q$ coefficient.
Thus $J+pb>A+qb$, proving (B-RD58.4), while
$K+b=A+qb$, proving (B-RD58.3). Its coefficient equation is
\[
 K f_K(z)=-(p-1)d\,d_U^{q-1}z.                           \tag{B-RD58.6}
\]
In particular $p\nmid K$ and this first coefficient is the same
nonzero scalar multiple of $z$ for every $z\ne0$.
No linearity of all the other coordinate coefficients has been
assumed.
\end{proof}

\begin{theorem}[Only residual projective order two remains]
\label{thm:B-RD58}
\status{Proved; full independent root and companion proof, original-source, and exact-code audits passed.}
In the normalized $p^2$ comparison with irreducible residual
inertia, different special characters require residual projective
order exactly two. In particular, if $p\equiv3\pmod4$, isomorphic
$p^2$-torsion in this irreducible residual case determines the
requested special character. If $p\equiv1\pmod4$, the only
remaining irreducible-residual type pair is the unique opposite
quadratic pair from Corollary~\ref{cor:B-RD53-quadratic}, and its
wild image must still be the full $M_2(\mathbf F_p)$.
At $p=5$, the pattern $\{A,B\}$ is excluded in the irreducible
residual case; only $\{A,C\}$ with wild order $625$ remains there.
\end{theorem}
\begin{proof}
Assume $h>2$. For a nonzero Cartan element of break $c>b$, the
degree-$K$ term from the preceding lemma changes before every
term of degree divisible by $p$:
\[
 \begin{split}
 J+pc-(K+c)&>(p-1)(c-b)>0,\\
 A+qc-(K+c)&=N(c-b)>0.
 \end{split}                                             \tag{B-RD58.7}
\]
It is also the first prime-to-$p$ term, so it gives the unique
first change. The required displacement valuation is $qA$.
Consequently $K+c=qA$ and the leading-coefficient equation on
the full final Cartan plane $Z$ is
\[
                   d_i^q\beta_i(g)z=K f_{i,K}(z)d_g.
                                                               \tag{B-RD58.8}
\]
Equation (B-RD58.6) yields constants $\kappa_i\ne0$ with
$\beta_i(g)=\kappa_i d_g$ for every $g\in Z$.
Since $Z$ is the full Cartan plane, $\beta_i(Z)=\mathbf F_q$.
But the different-type residual intertwiner is Frobenius-semilinear
by B53, and multiplication by $p$ uses the same orientation.
Thus $\beta_2(g)=\beta_1(g)^p$, just as in (B-RD57.4).
The two fixed scalar multiples of $d_g$ would make Frobenius
scalar-linear on all of $\mathbf F_q$, a contradiction.
Together with B57, this excludes every possible wild image at
$h>2$. Irreducibility excludes $h=1$.

For $h=2$, put $c_1\equiv1\pmod{p-1}$ for the first actual
tangent exponent. The residual eigenvalue is
$\lambda=\zeta^{c_1+p-1}$ and its projective order is the order
of $\lambda^{p-1}$. Order two means
\[
              (p-1)(c_1+p-1)\equiv N/2\pmod N.
\]
The different-type relation (B-RD53.4) then gives
$c_2-c_1\equiv N/2$, precisely an opposite quadratic pair.
Also $c_1+p-1$ is odd. For its class modulo $p+1$ to have
order two, $(p+1)/2$ must be odd, so $p\equiv1\pmod4$.
The cycle calculation of B53 gives exactly one such unordered
pair. The three-type table at $p=5$ gives the stated specialization.
\end{proof}

\begin{remark}[Scope and numbering]
\label{rem:B-RD58-scope}
\status{Proved scope and limitation; no pair-existence or least-index assertion.}
This argument uses upper quotient compatibility explicitly and only
then identifies the last lower group. It never sets a lower break
equal to the cyclotomic upper break without this passage.
Reducible residual representations remain outside its scope.
At $p=5$ the unrestricted least-index interval is still
$\nu_{\min}(5,4,2)\in\{1,2\}$, since the quadratic irreducible
case and the reducible residual case remain open.
The second lemma requires the irreducible tame action on the entire
anti-Cartan plane. Its scalar-$\mathbf F_q$ leading space is not
available automatically at projective order two, where one must
retain a general additive quadratic polynomial and possibly two
anti-Cartan breaks.
The retained exact supplement
\path{code/reduction_descent/verify_tame_full_wild_projective.py}
checks $24$ ramification-index cases, $5{,}574{,}772$ low-break
inequality cases, $1{,}180$ different-type cases, $92{,}456$
additive-polynomial coefficient cases, $364$ scalar/Frobenius cases,
$130{,}248$ primitive-point stabilizer cases and $3{,}026$ centralizer
cases. Its author run passed in $7.383174480994057$ seconds.
The complete independent root and companion audits passed; their
author-output-preserving exact reruns took $7.658215595998627$
and $7.812601506004285$ seconds, respectively.
These are finite supplements to the proof; no formal group,
division extension, or congruent pair is instantiated.
\end{remark}
